\section{The physical \texorpdfstring{\(Q^3\)}{Q3} column-energy gate}
\label{sec:tpc33-column-transfer}

The affine disintegration identifies the correct inner sequences.  We now
state exactly how much collective cancellation is needed from them.  The
outer row coefficient must not be included twice in the column energy.

\subsection{The two polarized column energies}

For a fixed right row \(\gamma\), define the left-ultra column
\begin{equation}
 B_\gamma^L
 :=\sum_\alpha\gamma_\alpha^{(1)}
 \sum_j\mathfrak A_{\alpha,\gamma}^{\rm act}(j)
 C_{m_\alpha}(j)\mathsf H_{m_\gamma}(j),
 \label{eq:tpc33-left-column}
\end{equation}
and its energy
\begin{equation}
 E_L:=\sum_\gamma|B_\gamma^L|^2.
 \label{eq:tpc33-left-column-energy}
\end{equation}
The coefficient \(\gamma_\gamma^{(2)}\) is outside
\(B_\gamma^L\).  Symmetrically, for a fixed left row \(\alpha\), put
\begin{align}
 B_\alpha^R
 &:=\sum_\gamma\gamma_\gamma^{(2)}
 \sum_j\mathfrak A_{\alpha,\gamma}^{\rm act}(j)
 \mathsf H_{m_\alpha}(j)C_{m_\gamma}(j),
 \label{eq:tpc33-right-column}\\
 E_R&:=\sum_\alpha|B_\alpha^R|^2.
 \label{eq:tpc33-right-column-energy}
\end{align}
Here \(\gamma_\alpha^{(1)}\) is outside \(B_\alpha^R\).  Define the two
scalar polarizations
\begin{equation}
 \cS_L=\sum_\gamma\gamma_\gamma^{(2)}B_\gamma^L,
 \qquad
 \cS_R=\sum_\alpha\gamma_\alpha^{(1)}B_\alpha^R.
 \label{eq:tpc33-scalar-polarizations}
\end{equation}
The exact identity \eqref{eq:tpc33-matched-polarization} gives
\begin{equation}
 \cS_{\rm sh}^{\rm all}=\cS_L+\cS_R.
 \label{eq:tpc33-shell-as-column-sum}
\end{equation}

\begin{theorem}[Matched polarized column-energy transfer]
\label{thm:tpc33-column-energy-transfer}
On the physical packet of \cref{sec:tpc33-physical-interface},
\begin{equation}
 \boxed{
 |\cS_{\rm sh}^{\rm all}|
 \ll_{\mathscr D}
 \sqrt{Q\log(2L)}
 \left(\sqrt{E_L}+\sqrt{E_R}\right).}
 \label{eq:tpc33-exact-energy-transfer}
\end{equation}
Consequently, the additional estimate
\begin{equation}
 \boxed{
 E_L+E_R\ll_{\eps,h,\mathscr D}X^\eps Q^3}
 \label{eq:tpc33-Q3-energy-input}
\end{equation}
for every fixed \(\eps>0\) implies
\begin{equation}
 \boxed{
 |\cS_{\rm sh}^{\rm all}|
 \ll_{\eps,h,\mathscr D}X^\eps Q^2.}
 \label{eq:tpc33-Q2-shell-output}
\end{equation}
\end{theorem}

\begin{proof}
By \eqref{eq:tpc33-scalar-polarizations}, Cauchy--Schwarz, and the inherited
row energy \eqref{eq:tpc33-row-l2},
\begin{align}
 |\cS_L|
 &\le
 \left(\sum_\gamma|\gamma_\gamma^{(2)}|^2\right)^{1/2}
 \left(\sum_\gamma|B_\gamma^L|^2\right)^{1/2}
 \notag\\
 &\ll_{\mathscr D}\sqrt{Q\log(2L)}\sqrt{E_L},
 \label{eq:tpc33-left-CS}
\end{align}
and similarly
\begin{equation}
 |\cS_R|
 \ll_{\mathscr D}\sqrt{Q\log(2L)}\sqrt{E_R}.
 \label{eq:tpc33-right-CS}
\end{equation}
Use \eqref{eq:tpc33-shell-as-column-sum} and the triangle inequality to
obtain \eqref{eq:tpc33-exact-energy-transfer}.  If
\eqref{eq:tpc33-Q3-energy-input} holds, then
\[
 \sqrt{E_L}+\sqrt{E_R}
 \le\sqrt{2(E_L+E_R)}
 \ll X^{\eps/2}Q^{3/2}.
\]
Applying the input with a smaller soft exponent and absorbing
\(\log(2L)\) proves \eqref{eq:tpc33-Q2-shell-output}.
\end{proof}

\begin{remark}[Sufficient, not equivalent]
The estimate \eqref{eq:tpc33-Q3-energy-input} controls the two
polarizations separately.  It is a clean sufficient condition for the
matched scalar shell, but cancellation between \(\cS_L\) and \(\cS_R\)
could make the scalar sum smaller even when one energy is large.  No
necessity or equivalence is asserted.
\end{remark}

\subsection{Splicing back to the distinguished zero coefficient}

Assume \(J\ll_{\mathscr D}Q\), as holds at the high-\(\beta\) endpoint in
\cref{sec:tpc33-high-beta-next-gate}.  Then \(C=\lfloor J\rfloor\) lies in the
inherited range of \eqref{eq:tpc33-zero-splice}.  Since \(X\asymp QJ\),
\begin{equation}
 Q^2+\frac{XQ}{J}\asymp Q^2.
 \label{eq:tpc33-content-at-J}
\end{equation}

\begin{corollary}[Conditional zero-frequency closure]
\label{cor:tpc33-zero-frequency-closure}
Suppose \(J\ll_{\mathscr D}Q\).  If
\eqref{eq:tpc33-Q3-energy-input} holds, then the small-content
distinguished coefficient from TPC-32 satisfies
\begin{equation}
 \boxed{
 |\widehat A_{\lfloor J\rfloor,q}(0)|
 \ll_{\eps,h,\mathscr D}X^\eps Q^2.}
 \label{eq:tpc33-zero-frequency-Q2}
\end{equation}
With the inherited sign convention, define the complete calibrated cutoff
difference by
\begin{equation}
 \cS_{\rm cal}^{\rm full}
 :=\cS_{\rm sh}^{\rm all}+\cD^{\rm dr}(\text{all}).
 \label{eq:tpc33-calibrated-full-definition}
\end{equation}
It satisfies
\begin{equation}
 \boxed{
 |\cS_{\rm cal}^{\rm full}|
 \ll_{\eps,h,\mathscr D}X^\eps
 \left(Q^2+\frac{XQ}{L}\right).}
 \label{eq:tpc33-calibrated-Q2}
\end{equation}
\end{corollary}

\begin{proof}
By \eqref{eq:tpc33-zero-splice},
\[
 |\widehat A_{\lfloor J\rfloor,q}(0)|
 \le |\cS_{\rm sh}^{\rm all}|
   +|\cS_{\rm sh}(c_{\alpha,\gamma}(j)>\lfloor J\rfloor)|.
\]
Use \eqref{eq:tpc33-Q2-shell-output},
\eqref{eq:tpc33-large-content-bound}, and
\eqref{eq:tpc33-content-at-J}.  The calibrated packet differs by the two
drift legs bounded in \eqref{eq:tpc33-drift-bound}; adding them proves
\eqref{eq:tpc33-calibrated-Q2}.
\end{proof}

This corollary is a direct amplitude estimate.  It matches the amplitude
budget associated with the \(1/400\) flatness gate in TPC-32, but it does
not prove the flatness-ratio inequality itself because no lower bound for
\(\sum_n|A_C(n)|^2\) is available.
